\documentclass[10pt]{extarticle}
\usepackage[a4paper,margin=1.5cm]{geometry}
\usepackage[T1]{fontenc}
\usepackage{amsmath,amssymb}
\usepackage{graphicx}
\usepackage[export]{adjustbox}
\usepackage{multicol}
\usepackage{enumitem}
\usepackage{booktabs,array,tabularx}
\setlength{\columnsep}{0.7cm}
\setlength{\parindent}{0pt}
\setlength{\parskip}{3pt}
\pagestyle{empty}
\makeatletter
% figure registry (figdefs.tex is generated by the extract step)
\newcommand{\qfigdef}[4]{\@namedef{qf@file@#1}{#2}\@namedef{qf@w@#1}{#3}\@namedef{qf@wide@#1}{#4}}
\newcommand{\qf@check}[1]{\@ifundefined{qf@file@#1}{\PackageError{score2}{Unknown figure id #1}{}}{}}
\newcommand{\qf@img}[4][]{\includegraphics[width=#2\dimexpr\@nameuse{qf@w@#4}pt\relax,max width=#3,max height=0.45\textheight,#1]{\@nameuse{qf@file@#4}}}
% Full-width questions: finalize wraps every question that contains a wide figure in \qfullbegin...\qfullend,
% so the WHOLE question (stem, figures, options) is set across the page in reading order, then the two columns
% resume. The question's head - its number, stem and first figure (\qfullsplit marks the end of the head) -
% is kept on one page, so a number is never left at a page foot with its figure overleaf; the rest may break
% across pages normally, so no big gaps are left. Without a \qfullsplit, the whole question is kept together
% if it is under 45% of a page. Never write these yourself.
\newif\ifqf@full
\newif\ifqf@boxed
\newsavebox{\qf@box}
\newcommand{\qf@open}{\global\qf@fulltrue\par\global\qf@boxedtrue
  \setbox\qf@box\vbox\bgroup\hsize\textwidth\linewidth\textwidth}
\newcommand{\qf@emit}[1]{\par\egroup\global\qf@boxedfalse
  \ifdim\dimexpr\ht\qf@box+\dp\qf@box\relax<#1\textheight\noindent\usebox{\qf@box}\par\else\unvbox\qf@box\fi}
\newcommand{\qfullsplit}{\ifqf@boxed\qf@emit{0.85}\fi}
\newcommand{\qf@close}{\ifqf@boxed\qf@emit{0.45}\fi\global\qf@fullfalse}
\newcommand{\qfullbegin}{\end{multicols}\qf@open}
\newcommand{\qfullend}{\qf@close\begin{multicols}{2}}
\newcommand{\qfullstart}{\qf@open}
\newcommand{\qfulldone}{\qf@close}
% \qfig[scale]{id}: block figure, centred, directly where it is written (size set by the tool).
% Fallback if a wide figure is not inside a full-width question: break out of the columns in place - no floats.
\newcommand{\qfig}[2][1]{\qf@check{#2}%
  \ifqf@full\qf@blk{#1}{#2}\else
  \ifnum\@nameuse{qf@wide@#2}=1 \expandafter\@firstoftwo\else\expandafter\@secondoftwo\fi
  {\end{multicols}\begin{center}\qf@img{#1}{\textwidth}{#2}\end{center}\begin{multicols}{2}}%
  {\qf@blk{#1}{#2}}\fi}
\newcommand{\qf@blk}[2]{\par\smallskip{\centering\qf@img{#1}{\linewidth}{#2}\par}\smallskip}
% \qfigi[scale]{id}: inline figure (inside option lists / tables / minipages); never breaks columns.
\newcommand{\qfigi}[2][1]{\qf@check{#2}\qf@img[valign=c]{#1}{\linewidth}{#2}}
\makeatother
% structure macros - lint reads these to check completeness and figure placement
\newcommand{\papersec}[1]{\par\bigskip{\centering\Large\bfseries PAPER-#1\par}\medskip}
\newcommand{\partsec}[2]{\par\medskip\noindent{\large\bfseries PART-#1: #2}\par\smallskip}
\newcommand{\sect}[1]{\par\smallskip\noindent\textbf{#1}\par}
\newcommand{\qn}[1]{\par\smallskip\noindent\textbf{#1.}~}
\newcommand{\paraq}[2]{\par\smallskip\noindent\textbf{Paragraph for Questions #1 and #2}\par}
\newenvironment{qoptions}{\begin{itemize}[leftmargin=1.8em,itemsep=1pt,topsep=2pt,label={}]}{\end{itemize}}
\newcommand{\opt}[2]{\item[(#1)] #2}
% figure registry (generated by score2_tools.py finalize)
\qfigdef{p1_phy_s1_q01_1}{Score Advanced - Test 02 (13 May 2021)-img/p1_phy_s1_q01_1.png}{105.5}{0}
\qfigdef{p1_phy_s1_q02_1}{Score Advanced - Test 02 (13 May 2021)-img/p1_phy_s1_q02_1.png}{107.2}{0}
\qfigdef{p1_phy_s1_q04_1}{Score Advanced - Test 02 (13 May 2021)-img/p1_phy_s1_q04_1.png}{56.6}{0}
\qfigdef{p1_phy_s1_q05_1}{Score Advanced - Test 02 (13 May 2021)-img/p1_phy_s1_q05_1.png}{105.8}{0}
\qfigdef{p1_phy_s2_q02_1}{Score Advanced - Test 02 (13 May 2021)-img/p1_phy_s2_q02_1.png}{103.0}{0}
\qfigdef{p1_phy_s2_q03_1}{Score Advanced - Test 02 (13 May 2021)-img/p1_phy_s2_q03_1.png}{99.9}{0}
\qfigdef{p1_phy_s2_q04_1}{Score Advanced - Test 02 (13 May 2021)-img/p1_phy_s2_q04_1.png}{103.8}{0}
\qfigdef{p1_phy_s2_q04_2}{Score Advanced - Test 02 (13 May 2021)-img/p1_phy_s2_q04_2.png}{105.5}{0}
\qfigdef{p1_phy_s2_q04_3}{Score Advanced - Test 02 (13 May 2021)-img/p1_phy_s2_q04_3.png}{104.1}{0}
\qfigdef{p1_phy_s2_q04_4}{Score Advanced - Test 02 (13 May 2021)-img/p1_phy_s2_q04_4.png}{106.6}{0}
\qfigdef{p1_phy_s3_q05_1}{Score Advanced - Test 02 (13 May 2021)-img/p1_phy_s3_q05_1.png}{101.3}{0}
\qfigdef{p1_phy_s3_q06_1}{Score Advanced - Test 02 (13 May 2021)-img/p1_phy_s3_q06_1.png}{131.1}{0}
\qfigdef{p1_phy_s3_q08_1}{Score Advanced - Test 02 (13 May 2021)-img/p1_phy_s3_q08_1.png}{158.4}{0}
\qfigdef{p1_chem_s1_q02_1}{Score Advanced - Test 02 (13 May 2021)-img/p1_chem_s1_q02_1.png}{291.5}{1}
\qfigdef{p1_chem_s1_q05_5}{Score Advanced - Test 02 (13 May 2021)-img/p1_chem_s1_q05_5.png}{40.9}{0}
\qfigdef{p1_chem_s1_q05_1}{Score Advanced - Test 02 (13 May 2021)-img/p1_chem_s1_q05_1.png}{38.2}{0}
\qfigdef{p1_chem_s1_q05_2}{Score Advanced - Test 02 (13 May 2021)-img/p1_chem_s1_q05_2.png}{38.9}{0}
\qfigdef{p1_chem_s1_q05_3}{Score Advanced - Test 02 (13 May 2021)-img/p1_chem_s1_q05_3.png}{38.5}{0}
\qfigdef{p1_chem_s1_q05_4}{Score Advanced - Test 02 (13 May 2021)-img/p1_chem_s1_q05_4.png}{65.9}{0}
\qfigdef{p1_chem_s2_q02_1}{Score Advanced - Test 02 (13 May 2021)-img/p1_chem_s2_q02_1.png}{107.2}{0}
\qfigdef{p1_chem_s2_q02_2}{Score Advanced - Test 02 (13 May 2021)-img/p1_chem_s2_q02_2.png}{103.1}{0}
\qfigdef{p1_chem_s2_q02_3}{Score Advanced - Test 02 (13 May 2021)-img/p1_chem_s2_q02_3.png}{107.2}{0}
\qfigdef{p1_chem_s2_q02_4}{Score Advanced - Test 02 (13 May 2021)-img/p1_chem_s2_q02_4.png}{104.1}{0}
\qfigdef{p1_chem_s3_q04_1}{Score Advanced - Test 02 (13 May 2021)-img/p1_chem_s3_q04_1.png}{185.2}{0}
\qfigdef{p1_chem_s3_q05_1}{Score Advanced - Test 02 (13 May 2021)-img/p1_chem_s3_q05_1.png}{308.0}{1}
\qfigdef{p1_chem_s3_q06_1}{Score Advanced - Test 02 (13 May 2021)-img/p1_chem_s3_q06_1.png}{86.3}{0}
\qfigdef{p2_phy_s1_q02_1}{Score Advanced - Test 02 (13 May 2021)-img/p2_phy_s1_q02_1.png}{102.1}{0}
\qfigdef{p2_phy_s1_q04_1}{Score Advanced - Test 02 (13 May 2021)-img/p2_phy_s1_q04_1.png}{53.6}{0}
\qfigdef{p2_phy_s1_q07_1}{Score Advanced - Test 02 (13 May 2021)-img/p2_phy_s1_q07_1.png}{128.9}{0}
\qfigdef{p2_phy_s2_para1_2_1}{Score Advanced - Test 02 (13 May 2021)-img/p2_phy_s2_para1_2_1.png}{130.3}{0}
\qfigdef{p2_phy_s3_q01_1}{Score Advanced - Test 02 (13 May 2021)-img/p2_phy_s3_q01_1.png}{88.7}{0}
\qfigdef{p2_phy_s3_q02_1}{Score Advanced - Test 02 (13 May 2021)-img/p2_phy_s3_q02_1.png}{184.9}{0}
\qfigdef{p2_phy_s3_q04_1}{Score Advanced - Test 02 (13 May 2021)-img/p2_phy_s3_q04_1.png}{83.2}{0}
\qfigdef{p2_phy_s3_q05_1}{Score Advanced - Test 02 (13 May 2021)-img/p2_phy_s3_q05_1.png}{157.3}{0}
\qfigdef{p2_phy_s3_q07_1}{Score Advanced - Test 02 (13 May 2021)-img/p2_phy_s3_q07_1.png}{95.1}{0}
\qfigdef{p2_phy_s3_q08_1}{Score Advanced - Test 02 (13 May 2021)-img/p2_phy_s3_q08_1.png}{91.1}{0}
\qfigdef{p2_chem_s1_q03_1}{Score Advanced - Test 02 (13 May 2021)-img/p2_chem_s1_q03_1.png}{262.5}{1}
\qfigdef{p2_chem_s1_q04_A}{Score Advanced - Test 02 (13 May 2021)-img/p2_chem_s1_q04_A.png}{60.5}{0}
\qfigdef{p2_chem_s1_q04_B}{Score Advanced - Test 02 (13 May 2021)-img/p2_chem_s1_q04_B.png}{102.1}{0}
\qfigdef{p2_chem_s1_q04_C}{Score Advanced - Test 02 (13 May 2021)-img/p2_chem_s1_q04_C.png}{107.0}{0}
\qfigdef{p2_chem_s1_q04_D}{Score Advanced - Test 02 (13 May 2021)-img/p2_chem_s1_q04_D.png}{100.7}{0}
\qfigdef{p2_chem_s1_q05_1}{Score Advanced - Test 02 (13 May 2021)-img/p2_chem_s1_q05_1.png}{309.3}{1}
\qfigdef{p2_chem_s1_q05_2}{Score Advanced - Test 02 (13 May 2021)-img/p2_chem_s1_q05_2.png}{62.0}{0}
\qfigdef{p2_chem_s1_q05_3}{Score Advanced - Test 02 (13 May 2021)-img/p2_chem_s1_q05_3.png}{31.4}{0}
\qfigdef{p2_chem_s3_q05_1}{Score Advanced - Test 02 (13 May 2021)-img/p2_chem_s3_q05_1.png}{168.2}{0}
\qfigdef{p2_chem_s3_q06_1}{Score Advanced - Test 02 (13 May 2021)-img/p2_chem_s3_q06_1.png}{195.7}{0}
\raggedcolumns
\makeatletter\@beginparpenalty=10000\makeatother
\newcolumntype{L}[1]{>{\raggedright\arraybackslash}p{#1}}
\begin{document}
\begin{center}{\Large\bfseries Score Advanced --- Test 02 (13 May 2021)}\end{center}
\begin{multicols}{2}
\raggedcolumns
\papersec{1}
\partsec{1}{PHYSICS}
\sect{SEC-1}
\qn{1} Two vertical strings having different lengths are joined to a uniform rod at points A and B. The rod is at rest. Now string connected to point B is burnt. Which of the following is/are correct just after this ?\par
\qfig{p1_phy_s1_q01_1}
\begin{qoptions}
\opt{A}{Acceleration of COM is horizontal.}
\opt{B}{Acceleration of COM is vertical.}
\opt{C}{Acceleration of particle at A is zero}
\opt{D}{Acceleration of particle at A is horizontal}
\end{qoptions}
\qn{2} A simple pendulum of length $\ell$ = 30 cm and mass m = 2kg is attached to a cart of mass 3 kg. The cart can move frictionlessly on level ground. The bob is released from rest when string makes an angle of 1.8$^\circ$ with the horizontal :-\par
\qfig{p1_phy_s1_q02_1}
\begin{qoptions}
\opt{A}{The amplitude of oscillation of cart is $\dfrac\pi5$ cm.}
\opt{B}{The time period of oscillation is $\dfrac\pi5$ sec}
\opt{C}{The cart and bob oscillate with phase difference of $\pi$.}
\opt{D}{The time period of oscillation is $\dfrac{\pi\sqrt3}5$ sec.}
\end{qoptions}
\qn{3} Consider the potential energy given by $\mathrm U=\dfrac{x^4}4-\dfrac{x^3}3-3x^2$
\begin{qoptions}
\opt{A}{x = 0 is position of stable equilibrium.}
\opt{B}{x = 3 is position of stable equilibrium.}
\opt{C}{If body starts from x = $-$2, it can go to x = 3 if its kinetic energy is greater than $\dfrac{16}3$ J.}
\opt{D}{If body starts from x = $-$2, it can go to x = 3 even if its initial kinetic energy is zero.}
\end{qoptions}
\qn{4} Four identical rods which have thermally insulated lateral surfaces are joined at point A. Points B,C,D \& E are connected to large reservoirs. If heat flows into the junction from point B at rate of 1 W and from point C at 3 W inside, flows out from D at 5 W, which relation(s) is/are correct for temperature of these points ?\par
\qfig{p1_phy_s1_q04_1}
\begin{qoptions}
\opt{A}{T$_\mathrm A<$ T$_\mathrm E$}
\opt{B}{T$_\mathrm B$ = T$_\mathrm C$}
\opt{C}{T$_\mathrm C>$ T$_\mathrm D$}
\opt{D}{T$_\mathrm B$ = T$_\mathrm E$}
\end{qoptions}
\qn{5} In a tank filled with water, a small hole is made on the side. The emerging water jet falls to a distance of 0.8 m from the base as shown (y is distance of hole from above) :-\par
\qfig{p1_phy_s1_q05_1}
\begin{qoptions}
\opt{A}{As the level of water falls, the range decreases.}
\opt{B}{As the level of water falls, the time of flight for water jet remains the same.}
\opt{C}{y can be 0.8 m.}
\opt{D}{y can be 0.2 m.}
\end{qoptions}
\qn{6} In a standard YDSE setup, the fringe width on the screen is 1.5 mm. When a thin glass film is pasted in front of the upper slit, the fringe pattern shifts up. But it is seen that at a point P above central maxima where intensity was one fourth the intensity at central maxima, intensity remains the same. There were no maxima between central maxima and point P before film was introduced. What can be the thickness of the film ? Take $\mu$ = 1.5, $\lambda$ = 450 nm :-
\begin{qoptions}
\opt{A}{3 $\times$ 10$^{-7}$m}
\opt{B}{5 $\times$ 10$^{-7}$ m}
\opt{C}{9 $\times$ 10$^{-7}$ m}
\opt{D}{1.4 $\times$ 10$^{-6}$ m}
\end{qoptions}
\sect{SEC-2}
\qn{1} A charged particle having non zero velocity is subjected to certain conditions given in List-I, List-II gives possible trajectories of the particle. Match the conditions in List-I with the results in List-II. Given all the fields are non-zero.\par
\noindent{\footnotesize\setlength{\tabcolsep}{2pt}\begin{tabular}{|L{.7cm}|L{3.0cm}|L{.7cm}|L{3.1cm}|}\hline
 & List-I & & List-II\\\hline
(P) & In only uniform electric field & (1) & The path of the charged particle may be a straight line\\\hline
(Q) & In only uniform magnetic field & (2) & The path of the charged particle may be a parabola\\\hline
(R) & In uniform magnetic and uniform electric field & (3) & The path of the charged particle may be circular\\\hline
(S) & Subjected to a net force of constant magnitude & (4) & The path of the charged particle may be a helical with uniform or non-uniform pitch.\\\hline
 & & (5) & None of these\\\hline
\end{tabular}}\par
\begin{qoptions}
\opt{A}{P $\to$ 1,2;Q $\to$ 1,3,4;R $\to$ 1,4;S $\to$ 1,2,3,4}
\opt{B}{P $\to$ 1,2,3;Q $\to$ 1,3,4;R $\to$ 1,4;S $\to$ 1,2,3}
\opt{C}{P $\to$ 1,2;Q $\to$ 1,2,3,4;R $\to$ 1,3,4;S $\to$ 1,2,3,4}
\opt{D}{P $\to$ 1,2,3;Q $\to$ 1,3;R $\to$ 1,3,4;S $\to$ 1,2,3}
\end{qoptions}
\qn{2} In the list--I, there are certain situations depicted and list--II gives the description of speed and acceleration of particle. Match them. Neglect air resistance and assume uniform gravitational field wherever applicable. Consider time before it hits the ground.\par
\noindent{\footnotesize\setlength{\tabcolsep}{2pt}\begin{tabular}{|L{.7cm}|L{3.4cm}|L{.7cm}|L{2.7cm}|}\hline
 & List-I & & List--II\\\hline
(P) & \qfigi{p1_phy_s2_q02_1}\par Particle moving in 1 dimension & (1) & speed increases \& then decreases\\\hline
(Q) & x = t$^2$ + 3t & (2) & speed decreases \& then increases\\\hline
(R) & A projectile thrown on level ground at an angle to the horizontal & (3) & speed increases continuously.\\\hline
(S) & A ball thrown vertically down from a cliff. & (4) & speed decreases continuously\\\hline
 & & (5) & acceleration is initially opposite to velocity.\\\hline
 & & (6) & acceleration is initially in the same direction as the velocity\\\hline
\end{tabular}}\par
\begin{qoptions}
\opt{A}{P $\to$ 1,5;Q $\to$ 2,6;R $\to$ 4;S $\to$ 2,5}
\opt{B}{P $\to$ 2,5;Q $\to$ 3,6;R $\to$ 2;S $\to$ 3,6}
\opt{C}{P $\to$ 1,5;Q $\to$ 4,6;R $\to$ 1,6;S $\to$ 2,5}
\opt{D}{P $\to$ 2,6;Q $\to$ 2;R $\to$ 3,6;S $\to$ 3,6}
\end{qoptions}
\qn{3} An ideal gas is taken along the reversible processes as represented by the adjoining diagram.\par
\qfig{p1_phy_s2_q03_1}
\noindent{\footnotesize\setlength{\tabcolsep}{2pt}\begin{tabular}{|L{.7cm}|L{3.4cm}|L{.7cm}|L{2.2cm}|}\hline
 & List-I & & List-II\\\hline
(P) & For process B $\to$ C & (1) & $\Delta$Q $>$ 0\\\hline
(Q) & For process A $\to$ B & (2) & $\Delta$W $>$ 0\\\hline
(R) & For cycle A $\to$ B $\to$ C $\to$ A & (3) & $\Delta$U $>$ 0\\\hline
(S) & For process C $\to$ A & (4) & $\Delta$W = 0\\\hline
 & & (5) & $\Delta$Q $<$ 0\\\hline
\end{tabular}}\par
\begin{qoptions}
\opt{A}{P $\to$ 1,2,3;Q $\to$ 1,2,5;R $\to$ 1,2,3;S $\to$ 1,2}
\opt{B}{P $\to$ 1,2;Q $\to$ 1,3,4;R $\to$ 1,4,5;S $\to$ 5}
\opt{C}{P $\to$ 1,2,3;Q $\to$ 1,3,4;R $\to$ 1,2;S $\to$ 5}
\opt{D}{P $\to$ 1,3;Q $\to$ 1,3,4;R $\to$ 5;S $\to$ 4,5}
\end{qoptions}
\qn{4} Match the list--I with nature of image in list--II (Consider all optical devices are kept in air)\par
\noindent{\footnotesize\setlength{\tabcolsep}{2pt}\begin{tabular}{|L{.7cm}|L{3.8cm}|L{.7cm}|L{1.8cm}|}\hline
 & List-I & & List-II\\\hline
(P) & \qfigi{p1_phy_s2_q04_1} & (1) & Erect image\\\hline
(Q) & \qfigi{p1_phy_s2_q04_2} & (2) & Inverted image\\\hline
(R) & \qfigi{p1_phy_s2_q04_3} & (3) & magnified image\\\hline
(S) & \qfigi{p1_phy_s2_q04_4} & (4) & diminished image\\\hline
 & & (5) & real image\\\hline
 & & (6) & virtual image\\\hline
\end{tabular}}\par
\begin{qoptions}
\opt{A}{P $\to$ 2,3,5;Q $\to$ 1,4,6;R $\to$ 1,3,6;S $\to$ 2,3,5}
\opt{B}{P $\to$ 1,4;Q $\to$ 2,3,5;R $\to$ 1,4,6;S $\to$ 1,3,6}
\opt{C}{P $\to$ 1,4,6;Q $\to$ 1,3,6;R $\to$ 2,3,5;S $\to$ 1,4}
\opt{D}{P $\to$ 2,3,5;Q $\to$ 1,3,6;R $\to$ 1,3,6;S $\to$ 1,4,6}
\end{qoptions}
\sect{SEC-3}
\qn{1} The side of a cube is measured by vernier callipers (10 divisions of a vernier scale coincide with 9 divisions of main scale, where 1 division of main scale is 1 mm). The main scale reads 10 mm and first division of vernier scale coincides with the main scale. Mass of the cube is 2.736 g. Find the density (in g/cm$^3$) of the cube in appropriate significant figures.
\qn{2} Three uniformly charged infinite wires with linear charge density $\lambda$ are placed along x, y and z axis respectively. Find the flux of electric field through Gaussian surface given by $x^2+y^2+z^2=1$ ; x $>$ 0 ; y $>$ 0 ; z $>$ 0. If your answer is $\dfrac{n\lambda}{12\epsilon_0}$ fill the value of n.
\qn{3} When a resistance R is connected in series with an element A, the electric current is found to be lagging behind the voltage by angle 30$^\circ$. When the same resistance is connected in series with element B, current leads by 60$^\circ$. When R, A, B are connected in series, the current now leads voltage by $\theta$ which is equal to $\tan^{-1}\left(\mathrm K/\sqrt3\right)$, then the value of K is (assume same AC source is used in all cases)
\qn{4} An x-ray beam of monochromatic photon are incident on a metallic surface having a negligible work function (take $\phi$ = 0). It is seen that the wavelength of most energetic photoelectrons is equal to the wavelength of x-ray photons. Find the wavelength (in pm). Round off to nearest integer.
\qn{5} Consider a circuit with an alternating source and contains inductor and capacitor. Given reading of A$_1$ and A$_2$ as 3 ampere and 5 ampere respectively. Find the reading of A in ampere.\par
\qfig{p1_phy_s3_q05_1}
\qn{6} In the circuit diagram shown, all the capacitors are in $\mu$F. The equivalent capacitance between points A and B is (in $\mu$ F).\par
\qfig{p1_phy_s3_q06_1}
\qn{7} A sample of $\beta$ active nuclei has a half life of 100 min. Initially, there are 10$^{13}$ $\beta$ active nuclei. Assume that all $\beta$ particles emitted leave the sample. What is the charge acquired (in $\mu$C) by the sample in 200 min?
\qn{8} A cube is immersed in kerosene and water at rest without any external force as shown. To keep it completely submerged in water, we need to apply an extra force of 4 N. If it was initially submerged to depth of $\dfrac a4$ in kerosene, what is the length `a' of the side of cube (in cm) ?\par
\qfig{p1_phy_s3_q08_1}
\partsec{2}{CHEMISTRY}
\sect{SEC-1}
\qn{1} Which statement(s) is/are correct -
\begin{qoptions}
\opt{A}{If 0.1 M CH$_3$COOH aqueous solution (K$_a$ = 1.8 $\times$ 10$^{-5}$) is diluted at 25$^\circ$C, then [H$^+$] will increase.}
\opt{B}{On heating a sample of pure water [H$^+$] and [OH$^-$] increases but the sample still remains neutral}
\opt{C}{10$^{-4}$ M HCl solution is less acidic than 0.1 M HCN (K$_{a\,\mathrm{HCN}}$ = 10$^{-5}$) solution}
\opt{D}{10$^{-8}$ M HCl aqueous solution will have pH exactly equal to 8}
\end{qoptions}
\qfullbegin
\qn{2} \qfig{p1_chem_s1_q02_1}
\qfullsplit
correct statements regarding the above reaction is/are :
\begin{qoptions}
\opt{A}{Free carbocation is involved in the reaction as an intermediate}
\opt{B}{Alkyl-oxygen bond cleavage takes place}
\opt{C}{Acyl-oxygen bond cleavage takes place}
\opt{D}{Alcohol formed in the reaction will give immediate turbidity with Lucas reagent}
\end{qoptions}
\qfullend
\qn{3} Which of the following reaction is/are feasible?
\begin{qoptions}
\opt{A}{2[Fe(CN)$_6$]$^{3-}$ + 2OH$^-$ + H$_2$O$_2$ $\to$ 2[Fe(CN)$_6$]$^{4-}$ + 2H$_2$O + O$_2$}
\opt{B}{2K$_4$[Fe(CN)$_6$] + H$_2$O$_2$ + H$_2$SO$_4$ $\to$ 2K$_3$[Fe(CN)$_6$] + K$_2$SO$_4$ + 2H$_2$O}
\opt{C}{2KMnO$_4$ + 3H$_2$O$_2$ $\to$ 2KOH + 2MnO$_2$ + 2H$_2$O + 3O$_2$}
\opt{D}{2KMnO$_4$ + 5H$_2$O$_2$ + 3H$_2$SO$_4$ $\to$ 2MnSO$_4$ + K$_2$SO$_4$ + 5O$_2$ + 8H$_2$O}
\end{qoptions}
\qn{4} Select the correct statement(s)?
\begin{qoptions}
\opt{A}{Addition of HgI$_2$ to aqueous solution of KI increases the freezing point since a complex K$_2$[HgI$_4$] is formed}
\opt{B}{Three aqueous solutions of ammonia labelled as A, B and C with concentrations 0.1 M, 0.01 M and 0.001 M respectively, then value of Van't Hoff factor for these solution will be in order :- i$_A$ $<$ i$_B$ $<$ i$_C$ (consider degree of dilution)}
\opt{C}{Vapour density of PCl$_5$(g) dissociating into PCl$_3$(g) and Cl$_2$(g) is 100. \% dissociation of PCl$_5$ is 4.25\% [Given : atomic mass of P = 31, Cl = 35.5]}
\opt{D}{If pH of 0.1 M monobasic acid is found to be 2, its osmotic pressure at a given temperature T K is 0.11 RT}
\end{qoptions}
\qn{5} Compounds having heat of hydrogenation lesser than that of cyclohexa-1, 4-diene i.e. \qfigi{p1_chem_s1_q05_5}\par
[Hint resonance energy of benzene is 36 Kcal]
\begin{qoptions}
\opt{A}{\qfigi{p1_chem_s1_q05_1}}
\opt{B}{\qfigi{p1_chem_s1_q05_2}}
\opt{C}{\qfigi{p1_chem_s1_q05_3}}
\opt{D}{\qfigi{p1_chem_s1_q05_4}}
\end{qoptions}
\qn{6} Which of the following compound/s does not liberate acidic gas during its hydrolysis.
\begin{qoptions}
\opt{A}{Ca$_3$P$_2$}
\opt{B}{AlN}
\opt{C}{Al$_2$S$_3$}
\opt{D}{CaH$_2$}
\end{qoptions}
\sect{SEC-2}
\qn{1} \noindent{\footnotesize\setlength{\tabcolsep}{2pt}\begin{tabular}{|L{.7cm}|L{3.0cm}|L{.7cm}|L{2.6cm}|}\hline
 & List-I (Electrolysis using inert electrode) & & List-II (Product at anode or cathode.)\\\hline
P & Aqueous solution of sodium acetate & 1 & H$_2$\\\hline
Q & Aqueous solution of sodium formate & 2 & CO$_2$\\\hline
R & Concentrated aqueous solution of NaCl & 3 & Cl$_2$\\\hline
S & Dilute aqueous solution of HCl & 4 & C$_2$H$_6$\\\hline
 & & 5 & O$_2$\\\hline
\end{tabular}}\par
The correct option is -
\begin{qoptions}
\opt{A}{P $\to$ 1,2,4;Q $\to$ 1,2;R $\to$ 1, 3;S $\to$ 1, 5}
\opt{B}{P $\to$ 1,5;Q $\to$ 2,4;R $\to$ 1;S $\to$ 3, 4}
\opt{C}{P $\to$ 1,4;Q $\to$ 2,4;R $\to$ 1, 5;S $\to$ 3, 5}
\opt{D}{P $\to$ 2,4;Q $\to$ 1,4;R $\to$ 2, 5;S $\to$ 1,5}
\end{qoptions}
\qn{2} Match the list\par
\noindent{\footnotesize\setlength{\tabcolsep}{2pt}\begin{tabular}{|L{.7cm}|L{3.3cm}|L{.7cm}|L{2.2cm}|}\hline
 & List-I & & List-II\\\hline
P & \qfigi{p1_chem_s2_q02_1} & 1 & Ionic organometalic product is formed\\\hline
Q & \qfigi{p1_chem_s2_q02_2} & 2 & Product can show electrophillic addition reaction with HBr (1 eq.)\\\hline
R & \qfigi{p1_chem_s2_q02_3} & 3 & Reaction involve free radical as an intermediate\\\hline
S & \qfigi{p1_chem_s2_q02_4} & 4 & Number of carbon in parent chain of product is 4\\\hline
\end{tabular}}\par
Select correct code for your answer :
\begin{qoptions}
\opt{A}{P $\to$ 3, 4;Q $\to$ 2;R $\to$ 2, 4;S $\to$ 1, 4}
\opt{B}{P $\to$ 3, 4;Q $\to$ 2;R $\to$ 2, 4;S $\to$ 1, 3}
\opt{C}{P $\to$ 3, 4;Q $\to$ 1, 2;R $\to$ 2, 4;S $\to$ 1, 4}
\opt{D}{P $\to$ 2, 3;Q $\to$ 2;R $\to$ 2, 4;S $\to$ 1, 4}
\end{qoptions}
\qn{3} Match the list.\par
\noindent{\footnotesize\setlength{\tabcolsep}{2pt}\begin{tabular}{|L{.7cm}|L{2.9cm}|L{.7cm}|L{2.9cm}|}\hline
 & List-I (complex) & & List-II (Properties)\\\hline
P & [CoF$_6$]$^{3-}$ & 1 & Diamagnetic and low spin complex\\\hline
Q & [Cr(C$_2$O$_4$)$_3$]$^{3-}$ & 2 & Paramagnetic and outer orbital complex\\\hline
R & [AuCl$_4$]$^-$ & 3 & Paramagnetic and inner orbital complex\\\hline
S & [Fe(CN)$_6$]$^{3-}$ & 4 & Paramagnetic and low spin complex\\\hline
\end{tabular}}\par
The correct option is
\begin{qoptions}
\opt{A}{P $\to$ 1;Q $\to$ 2;R $\to$ 3;S $\to$ 4}
\opt{B}{P $\to$ 2;Q $\to$ 3;R $\to$ 1;S $\to$ 4}
\opt{C}{P $\to$ 2;Q $\to$ 1;R $\to$ 3;S $\to$ 4}
\opt{D}{P $\to$ 1;Q $\to$ 3;R $\to$ 2;S $\to$ 4}
\end{qoptions}
\qn{4} Match the list:\par
\noindent{\footnotesize\setlength{\tabcolsep}{2pt}\begin{tabular}{|L{.7cm}|L{3.1cm}|L{.7cm}|L{2.7cm}|}\hline
 & List-I (Complex Compound ) & & List- II (Properties)\\\hline
P & [Ni(H$_2$O)$_6$]Cl$_2$ & 1 & d$^2$sp$^3$ hybridisation\\\hline
Q & [Co(CN)$_2$(NH$_3$)$_4$] OC$_2$H$_5$ & 2 & Ionisation isomerism\\\hline
R & [IrCl$_6$]$^{3-}$ & 3 & $\mu$ = 2.83 B.M.\\\hline
S & [PtCl$_2$(NH$_3$)$_4$]Br$_2$ & 4 & $\Delta_0$ $<$ P\\\hline
 & & 5 & sp$^3$d$^2$ hybridisation\\\hline
\end{tabular}}\par
The correct option is
\begin{qoptions}
\opt{A}{P $\to$ 3,4,5;Q $\to$ 1,2;R $\to$ 1;S $\to$ 1,2}
\opt{B}{P $\to$ 3,4;Q $\to$ 1;R $\to$ 2,3;S $\to$ 1,2}
\opt{C}{P $\to$ 1,2,3;Q $\to$ 1,4;R $\to$ 4,5;S $\to$ 2,5}
\opt{D}{P $\to$ 2,3;Q $\to$ 2;R $\to$ 3,4;S $\to$ 1}
\end{qoptions}
\sect{SEC-3}
\qn{1} For the first order reaction A(g) $\to$ 2B(g) + C(g) the half life for decomposition of A is 3 minute at 300K. Calculate the time (in minute) in which partial pressure of A(g) will drop from 2 bar to 0.5 bar at 400K. Given activation energy of reaction is 840 R. (Take ln2 = 0.7)
\qn{2} 2 moles of methane gas are compressed from (1 bar, 2L) to 2 bar isothermally against constant external pressure. Calculate work done in Joules.
\qn{3} For 10 minutes each, at 0$^\circ$C, from two identical holes nitrogen and an unknown gas are leaked into a common vessel of 4 litre capacity. The resulting pressure in vessel is 2.8 atm and the mixture contains 0.4 mole of nitrogen. What is the molar mass of unknown gas? (Take R = 0.0821 L--atm - mol$^{-1}$ K$^{-1}$)
\qn{4} \qfig{p1_chem_s3_q04_1}
Find out total number of final products?
\qfullbegin
\qn{5} How many of the following compounds given below have faster aromatic nucleophilic substitution\par
\qfig{p1_chem_s3_q05_1}
\qfullend
\qn{6} Find out total number of moles of Grignard reagent RMgX consumed when treated it with 1 mole of `X'\par
\qfig{p1_chem_s3_q06_1}
\qn{7} Find the number of electrons having $|$m$|$ $\leq$ 1, l = 2, m$_s$ = $-\dfrac{1}{2}$ in Pd (Z = 46).
\qn{8} During extraction of iron, CO is the chief reducing agent. If 72 Kg of carbon is available and produced CO is completely used to reduce Haematite (present in excess amount) then calculate how much Kg of Iron is obtained after reduction.\par
(Given : atomic weight of C--12 g/mol, O--16 g/mol, Fe--55.85 g/mol)
\partsec{3}{MATHEMATICS}
\sect{SEC-1}
\qn{1} If e$_1$ is eccentricity of 9x$^2$ + 4y$^2$ = 36 and e$_2$ is eccentricity of 9x$^2$ $-$ 4y$^2$ = 36 then
\begin{qoptions}
\opt{A}{e$_1^2$ + e$_2^2$ $<$ 3}
\opt{B}{e$_1^2$ + e$_2^2$ $>$ 3}
\opt{C}{e$_1^2$ + e$_2^2$ $<$ 4}
\opt{D}{e$_1^2$ + e$_2^2$ $>$ 4}
\end{qoptions}
\qn{2} Let $f : [0, 1] \to [0, \infty)$ be a differentiable function so that $f'(x)$ is decreasing function ; $f(0) = 0$, $f'(1) > 0$, $f'(0) > 0$ and $I = \displaystyle\int_0^1 \dfrac{dx}{1 + f^2(x)}\,dx$ then which of the following is INCORRECT ?
\begin{qoptions}
\opt{A}{$I \leqslant \dfrac{f(1)}{f'(1)}$}
\opt{B}{$I < \dfrac{f(1)}{f'(1)}$}
\opt{C}{$I \geqslant \dfrac{f(1)}{f'(1)}$}
\opt{D}{$I > \dfrac{f(1)}{f'(1)}$}
\end{qoptions}
\qn{3} The coordinates of four angular points of a tetrahedron OABC are (0,0,0), (0,0,2), (0,4,0) and (6,0,0) respectively. A point P inside the tetrahedron is at the same distance r from the four plane faces of tetrahedron. Which of following cannot be the value of r ?
\begin{qoptions}
\opt{A}{$\dfrac{2}{3}$}
\opt{B}{$\dfrac{2}{5}$}
\opt{C}{2}
\opt{D}{3}
\end{qoptions}
\qn{4} Let C$_k$ = $^n$C$_k$ for 0 $\leq$ k $\leq$ n and A$_k$ = $\begin{bmatrix} C_{k-1} & 0\\ 0 & C_k \end{bmatrix}$ for k $\geq$ 1 and\par
A$_1$A$_2$ + A$_2$A$_3$ + A$_3$A$_4$ + ..... + A$_{n-1}$A$_n$ = $\begin{bmatrix} a & 0\\ 0 & b \end{bmatrix}$ then
\begin{qoptions}
\opt{A}{a = b}
\opt{B}{a + b = 2 ($^{2n}$C$_{n-1}$ $-$ n)}
\opt{C}{a = 2b}
\opt{D}{a + b = 0}
\end{qoptions}
\qn{5} Probability that the length of a randomly chosen chord of a circle lies between $\dfrac{1}{2}$ and $\dfrac{3}{4}$ of its diameter is $\dfrac{p}{q}$ where p and q are coprime number, then
\begin{qoptions}
\opt{A}{p + q = 21}
\opt{B}{q $-$ p = 11}
\opt{C}{p + q = 20}
\opt{D}{q $-$ p = 10}
\end{qoptions}
\qn{6} For two complex numbers z$_1$ and z$_2$ ; (az$_1$ + b$\bar{z}_1$) (cz$_2$ + d$\bar{z}_2$) = (cz$_1$ + d$\bar{z}_1$) (az$_2$ + b$\bar{z}_2$) if
\begin{qoptions}
\opt{A}{$\dfrac{a}{b} = \dfrac{c}{d}$}
\opt{B}{$\dfrac{a}{d} = \dfrac{b}{c}$}
\opt{C}{$|z_1| = |z_2|$}
\opt{D}{arg(z$_1$) = arg(z$_2$)}
\end{qoptions}
\sect{SEC-2}
\qn{1} Let $\phi(x) = 3f\left(\dfrac{x^2}{3}\right) + f\left(3 - x^2\right)\ \forall\ x \in (-3, 4)$ where $f''(x) > 0\ \forall\ x \in (-3, 4)$,\par
$g(x) = (a^2 - 3a + 2)\left(\cos^2\dfrac{x}{4} - \sin^2\dfrac{x}{4}\right) + (a - 1)x + \sin 1$\par
Match List-I with List-II and select the correct answer using the code given below the list.\par
\noindent{\footnotesize\setlength{\tabcolsep}{2pt}\begin{tabular}{|L{.7cm}|L{2.6cm}|L{.7cm}|L{2.9cm}|}\hline
(P) & $\phi(x)$ is increasing in & (1) & $(0, 1) \cup (1, 4)$\\\hline
(Q) & $\phi(x)$ is decreasing in & (2) & $\left[-\dfrac{3}{2}, 0\right]\ ;\ \left[\dfrac{3}{2}, 4\right)$\\\hline
(R) & The set of values of a for which g(x) does not possess critical point is & (3) & $\left(-3, -\dfrac{3}{2}\right]\ ;\ \left[0, \dfrac{3}{2}\right]$\\\hline
(S) & The set of values of a for which g(x) possess critical point is & (4) & $(-\infty, 0] \cup \{1\} \cup [4, \infty)$\\\hline
\end{tabular}}\par
\begin{qoptions}
\opt{A}{P $\to$ 2;Q $\to$ 3;R $\to$ 1;S $\to$ 4}
\opt{B}{P $\to$ 2;Q $\to$ 3;R $\to$ 4;S $\to$ 1}
\opt{C}{P $\to$ 3;Q $\to$ 2;R $\to$ 1;S $\to$ 4}
\opt{D}{P $\to$ 3;Q $\to$ 2;R $\to$ 4;S $\to$ 1}
\end{qoptions}
\qn{2} S $\equiv$ x$^2$ + y$^2$ $-$ 2x $-$ 4y + 1 = 0 and S$'$ $\equiv$ x$^2$ + y$^2$ $-$ 4x $-$ 2y + 1 = 0 are equations of two circles.\par
Match List-I with List-II and select the correct answer using the code given below the list.\par
\noindent{\footnotesize\setlength{\tabcolsep}{2pt}\begin{tabular}{|L{.7cm}|L{3.8cm}|L{.7cm}|L{1.5cm}|}\hline
(P) & Sum of number of tangents from (1, 3) to S and number of tangents from (1,3) to S$'$ is & (1) & 1\\\hline
(Q) & If r is radius of circle passes through point of intersection of S and S$'$ and whose normal is the straight line 2x + y $-$ 6 = 0, then the value of $\dfrac{r^2}{4}$ is & (2) & 2\\\hline
(R) & If m is slope of common tangent of S = 0 and S$'$ = 0 then $|$m$|$ is & (3) & 3\\\hline
(S) & If $\ell$ is length of common tangent of S and S$'$ then the value of 2$\ell^2$ is & (4) & 4\\\hline
\end{tabular}}\par
\begin{qoptions}
\opt{A}{P $\to$ 2;Q $\to$ 1;R $\to$ 2;S $\to$ 3}
\opt{B}{P $\to$ 2;Q $\to$ 2;R $\to$ 1;S $\to$ 4}
\opt{C}{P $\to$ 2;Q $\to$ 2;R $\to$ 4;S $\to$ 1}
\opt{D}{P $\to$ 1;Q $\to$ 1;R $\to$ 4;S $\to$ 2}
\end{qoptions}
\qn{3} (1) The sum of all the values of r satisfying $^{39}$C$_{3r-1}$ $-$ $^{39}$C$_{r^2}$ = $^{39}$C$_{r^2-1}$ $-$ $^{39}$C$_{3r}$ is $\alpha_1$.\par
(2) If $^{2n+3}$C$_{2n}$ $-$ $^{2n+2}$C$_{2n-1}$ = 15(2n + 1) then the value of n is $\alpha_2$\par
(3) If $^{56}$P$_{r+6}$ : $^{54}$P$_{r+3}$ = 30800 : 1 then the value of r is $\alpha_3$\par
(4) $^{n+2}$C$_8$ : $^{n-2}$P$_4$ = 57 : 16 then the value of n is $\alpha_4$\par
Match List-I with List-II and select the correct answer using the code given below the list.\par
\noindent{\footnotesize\setlength{\tabcolsep}{2pt}\begin{tabular}{|L{.7cm}|L{3.2cm}|L{.7cm}|L{1.5cm}|}\hline
(P) & The value of $\alpha_1$ is & (1) & 41\\\hline
(Q) & The value of $\alpha_2$ is & (2) & 8\\\hline
(R) & The value of $\alpha_3$ is & (3) & 14\\\hline
(S) & The value of $\alpha_4$ is & (4) & 19\\\hline
\end{tabular}}\par
\begin{qoptions}
\opt{A}{P $\to$ 2;Q $\to$ 3;R $\to$ 1;S $\to$ 4}
\opt{B}{P $\to$ 2;Q $\to$ 3;R $\to$ 4;S $\to$ 1}
\opt{C}{P $\to$ 2;Q $\to$ 1;R $\to$ 3;S $\to$ 4}
\opt{D}{P $\to$ 3;Q $\to$ 2;R $\to$ 1;S $\to$ 4}
\end{qoptions}
\qn{4} (i) Let harmonic mean, arithmetic mean and geometric mean of two positive numbers a and b respectively be 4, A and G. If 2A + G$^2$ = 27 then a$^2$ + b$^2$ = $\alpha_1$\par
(ii) If the value of x + y + z = 15 where a, x, y, z, b are in arithmetic progression while the value of $\dfrac{1}{x} + \dfrac{1}{y} + \dfrac{1}{z}$ is $\dfrac{5}{8}$ where a, x, y, z, b are in harmonic progression, then the value of a$^2$ + b$^2$ is $\alpha_2$\par
(iii) The value of $\displaystyle\sum_{r=1}^{16}\left(\dfrac{1^3 + 2^3 + \ldots + r^3}{1 + 3 + \ldots + (2r - 1)}\right)$ is $\alpha_3$\par
(iv) If the geometric mean and harmonic mean of two positive numbers x$_1$ and x$_2$ are 18 and $\dfrac{216}{13}$ respectively then the value of $|x_1 - x_2|$ is $\alpha_4$.\par
Match List-I with List-II and select the correct answer using the code given below the list.\par
\noindent{\footnotesize\setlength{\tabcolsep}{2pt}\begin{tabular}{|L{.7cm}|L{3.2cm}|L{.7cm}|L{1.5cm}|}\hline
(P) & The value of $\alpha_1$ is & (1) & 52\\\hline
(Q) & The value of $\alpha_2$ is & (2) & 45\\\hline
(R) & The value of $\alpha_3$ is & (3) & 446\\\hline
(S) & The value of $\alpha_4$ is & (4) & 15\\\hline
\end{tabular}}\par
\begin{qoptions}
\opt{A}{P $\to$ 2;Q $\to$ 1;R $\to$ 2;S $\to$ 3}
\opt{B}{P $\to$ 2;Q $\to$ 2;R $\to$ 1;S $\to$ 4}
\opt{C}{P $\to$ 2;Q $\to$ 1;R $\to$ 3;S $\to$ 4}
\opt{D}{P $\to$ 1;Q $\to$ 1;R $\to$ 4;S $\to$ 2}
\end{qoptions}
\sect{SEC-3}
\qn{1} In a triangle ABC (with usual notations), a = 5, b = 4 and A $-$ B = $\cos^{-1}\left(\dfrac{31}{32}\right)$ then the value of $\dfrac{c}{5}$ is
\qn{2} In the expansion of (1 $-$ x)$^{30}$(1 + x + x$^2$)$^{31}$ the coefficient of x$^{47}$ is ($-^{30}$C$_r$) then the value of $\left|\dfrac{r}{2}\right|$ is
\qn{3} In the equation x$^5$ $-$ 5x$^4$ + 9x$^3$ $-$ 9x$^2$ + 5x $-$ 1 = 0 if p is the sum of real roots and q is the product of real roots then the value of $\dfrac{q}{p}$ is
\qn{4} Let $f : (0, \infty) \to (0, \infty)$ be a function differentiable at 3 and satisfying $f(3) = 3f'(3)$. If the value of $\displaystyle\lim_{x \to \infty}\left(\dfrac{f\left(3 + \frac{3}{x}\right)}{f(3)}\right)^x$ is L then $\dfrac{[L]}{5}$ is (where [.] is greatest integer function)
\qn{5} The value of $\displaystyle\int_{\pi/4}^{3\pi/4}\dfrac{\cos x}{1 - \cos x}\,dx$ is equal to (use $\pi = \dfrac{22}{7}$)
\qn{6} Let $f(x) =$\par{\scriptsize$\begin{vmatrix} \sec x & \cos x & \sec^2 x + \cot x\,\mathrm{cosec}\,x + \cos x\\ \cos^2 x & \cos^2 x & \mathrm{cosec}^2 x + \cos^4 x\\ 1 & \cos^2 x & \cos^2 x \end{vmatrix}$}\par then the value $\left|\displaystyle\int_0^{\pi/2} f(x)\,dx\right|$ is (use $\pi = \dfrac{22}{7}$)
\qn{7} If y = 2x $-$ 3 is a tangent to the parabola y$^2$ = 4a $\left(x - \dfrac{1}{3}\right)$ then the value of $|$a$|$ is
\qn{8} OA, OB, OC are the sides of a rectangular parallelopiped whose diagonals are OO$'$, AA$'$, BB$'$ and CC$'$. D is the centre of rectangle AC$'$O$'$B and D$'$ is centre of rectangle O$'$A$'$CB$'$. If sides OA, OB, OC are in ratio 1 : 2 : 3 and $\angle$DOD$'$ = $\alpha$ then $\dfrac{697\cos^2\alpha}{128}$ is\par
(Centre of rectangle is mid point of diagonal)
\raggedcolumns
\papersec{2}
\partsec{1}{PHYSICS}
\sect{SEC-1}
\qn{1} Pick the correct option/options.
\begin{qoptions}
\opt{A}{If the sum of all the forces acting on a particle is zero, then it will necessarily be in equilibrium.}
\opt{B}{When a sphere rolls on a horizontal fixed surface, there is no point on the sphere which has vertical velocity.}
\opt{C}{When tall buildings are constructed on earth, the duration of day and night increases.}
\opt{D}{A body is rotating non uniformly about a vertical axis fixed in an inertial reference frame. The resultant force on the particle of the body not on the axis is horizontal and skew with the axis.}
\end{qoptions}
\qn{2} A charged ball of charge +q$_1$ is revolving around another charge +q as shown in a conical pendulum. The motion is in a horizontal plane :-\par
\qfig{p2_phy_s1_q02_1}
\begin{qoptions}
\opt{A}{Tension in the string is greater than the weight of the ball.}
\opt{B}{The tension in the string is greater than the electrostatic repulsive force.}
\opt{C}{If the charge is removed, the speed of ball has to be increased to maintain the angle.}
\opt{D}{If the charge is removed, the speed of ball has to be decreased to maintain the angle.}
\end{qoptions}
\qn{3} Mark the incorrect statement(s) for a standing wave established in a string :
\begin{qoptions}
\opt{A}{Energy of an element of medium remains constant.}
\opt{B}{Phase difference between any two particles can be $\dfrac{\pi}{2}$}
\opt{C}{Energy can get transferred through antinodes.}
\opt{D}{Energy can get transferred through nodes.}
\end{qoptions}
\qn{4} A wall is moving with velocity u and a source of sound (frequency f ) moves with velocity $\dfrac{u}{2}$ in the same direction as shown. Assuming that the sound travels with a speed of 10u, pick the right option or options.\par
\qfig{p2_phy_s1_q04_1}
\begin{qoptions}
\opt{A}{wavelength of incident sound is $\dfrac{19u}{2f}$.}
\opt{B}{frequency of the incident sound is $\dfrac{18}{19}$f.}
\opt{C}{wavelength of reflected sound is $\dfrac{u}{f}$.}
\opt{D}{frequency of the reflected sound is $\dfrac{12f}{19}$.}
\end{qoptions}
\qn{5} A gas is undergoing an adiabatic process. At a certain stage, the values of volume and temperature are V$_0$ and T$_0$ respectively and the modulus of slope of V--T curve is m. Choose the correct options-
\begin{qoptions}
\opt{A}{The value of C$_V$ = $\dfrac{RT_0 m}{V_0}$}
\opt{B}{The value of C$_P$ = $\left(1 + \dfrac{T_0 m}{V_0}\right)$R}
\opt{C}{m = $\dfrac{V_0}{T_0(\gamma - 1)}$}
\opt{D}{m = $\dfrac{-V_0}{T_0(\gamma - 1)}$}
\end{qoptions}
\qn{6} A long straight cylindrical shell has inner radius R$_i$ and outer radius R$_0$. It carries current i, uniformly distributed over its cross-section. A solid wire of radius r is to be placed parallel to the cylinder axis, in the hollow region such that (r $<$ R$_i$) :-
\begin{qoptions}
\opt{A}{The magnetic field is zero everywhere in hollow region. We conclude that the wire is on the cylinder axis and carries current i in the same direction as current in the shell.}
\opt{B}{The magnetic field is zero everywhere outside the shell (r $>$ R$_0$). We conclude that the wire is on the cylinder axis and carries i in the direction opposite to the direction of current in the shell.}
\opt{C}{The magnetic field is zero everywhere in the hollow region. We conclude that the wire may be anywhere in the hollow region but must be carrying current i in the same direction as the current in the shell.}
\opt{D}{The magnetic field is zero everywhere in the hollow region. We conclude that wire does not carry any current.}
\end{qoptions}
\qn{7} The diagram shows a modified meter bridge, which is used for measuring two unknown resistance at the same time. When only the first galvanometer P is used, for obtaining the balance point, it is found at point C. Now the first galvanometer P is removed and the second galvanometer Q is used, which gives balance point at D. Using the details given in the diagram, find out the value of R$_1$ and R$_2$.\par
\qfig{p2_phy_s1_q07_1}
\begin{qoptions}
\opt{A}{R$_1$ = 5R/3}
\opt{B}{R$_2$= 4R/3}
\opt{C}{R$_1$ = 4R/3}
\opt{D}{R$_2$ = 5R/3}
\end{qoptions}
\qn{8} A mirror is hanging by a light string of length $\ell$ from ceiling. A light pulse of monochromatic light rays is incident normally on the mirror and gets reflected. Which of the following statements is/are correct ?
\begin{qoptions}
\opt{A}{Mirror will start to oscillate.}
\opt{B}{Wavelength of reflected rays will be greater than that of the incident rays.}
\opt{C}{Wavelength of reflected rays may be less than that of incident rays.}
\opt{D}{Assuming the intensity of incident light to be constant, the maximum angle to which the mirror deflects depends on the wavelength of the light pulse.}
\end{qoptions}
\sect{SEC-2}
\paraq{1}{2}For given circuit, give answer of following questions\par
\qfig{p2_phy_s2_para1_2_1}
\qn{1} If device--1 is battery of emf 10V and 2$\Omega$ internal resistance (+ive end at A) and device-2 is battery of emf x volt and internal resistance 1$\Omega$ (+ive end at C) and reading of galvanometer is zero then value of x is :-
\begin{qoptions}
\opt{A}{6 volt}
\opt{B}{4 volt}
\opt{C}{10 volt}
\opt{D}{2 volt}
\end{qoptions}
\qn{2} If device-1 and device-2 are capacitors of capacitance C$_1$ and C$_2$ and reading of galvanometer remain zero then $\dfrac{C_1}{C_2}$ is (no charge pass through Galvanometer):-
\begin{qoptions}
\opt{A}{$\dfrac{1}{3}$}
\opt{B}{$\dfrac{1}{2}$}
\opt{C}{1}
\opt{D}{$\dfrac{2}{3}$}
\end{qoptions}
\paraq{3}{4}A charged particle having mass `m' and charge `+q' is released from rest at origin in uniform electric field $\vec{E} = E_0\hat{j}$ and uniform magnetic field $\vec{B} = B_0\hat{k}$.\par
\qn{3} Maximum value of `y' coordinate of particle in the path.
\begin{qoptions}
\opt{A}{$\dfrac{mE_0}{qB_0^2}$}
\opt{B}{2$\dfrac{mE_0}{qB_0^2}$}
\opt{C}{$\dfrac{3mE_0}{qB_0^2}$}
\opt{D}{$\dfrac{4mE_0}{qB_0^2}$}
\end{qoptions}
\qn{4} Mark the correct option(s) :
\begin{qoptions}
\opt{A}{When y-coordinate is maximum, the speed of the particle is maximum.}
\opt{B}{When y-coordinate is maximum, the velocity of particle is in x-direction.}
\opt{C}{Radius of curvature of path when y-coordinate is maximum is $\dfrac{4mE_0}{qB_0^2}$.}
\opt{D}{Radius of curvature of path when y-coordinate is maximum is $\dfrac{2mE_0}{qB_0^2}$.}
\end{qoptions}
\sect{SEC-3}
\qn{1} A uniform solid hemisphere of mass m and radius R is attached to the roof with a chord of torsional constant C and performing torsional SHM. Then the time period (in seconds) of SHM is (Take m = 15 kg, R = $\dfrac{2}{\pi}$m, C = 6 Nm/rad).\par
\qfig{p2_phy_s3_q01_1}
\qn{2} Young's double slit experiment is conducted in a liquid of refractive index $\mu_1$ as shown in figure. A thin transparent slab of refractive index $\mu_2$ is placed in front of the slit S$_2$. If $\mu_2 = \dfrac{3}{2}$, $\mu_1 = \dfrac{4}{3}$, t = 8 $\mu$m, the magnitude of optical path difference (wrt liquid) at `O' in $\mu$m is :\par
\qfig{p2_phy_s3_q02_1}
\qn{3} We would like to increase the length of a 15 cm long copper rod of cross-section 4 mm$^2$ by 1mm. The energy absorbed by the rod if it is heated is E$_1$. The energy absorbed by the rod if it is stretched slowly is E$_2$. The ratio of E$_1$ and E$_2$ is $\alpha$. The value of $\dfrac{\alpha}{100}$ is. [Various parameters of Copper are: Density = 9 $\times$ 10$^3$ kg/m$^3$, Thermal coefficient of linear expansion = 16 $\times$ 10$^{-6}$ K$^{-1}$, Young's modulus = 135 $\times$ 10$^9$ Pa, Specific heat = 400 J/kg-K]
\qn{4} A pendulum bob hinged at O is given velocity u = $\sqrt{4g\ell}$ from bottommost point as shown. If $\ell$ = 54 cm, find maximum height (in m) reached by bob from point of projection.\par
\qfig{p2_phy_s3_q04_1}
\qn{5} A concave mirror of focal length 20 cm and a convex lens of focal length 10 cm are kept with their optic axis parallel but separated by 0.5 mm as shown in figure. The distance between lens and mirror is 10 cm. An object of height 3 mm is placed on the optic axis of lens at distance 15 cm from the lens. Find height of image formed by mirror after refraction from lens (in mm).\par
\qfig{p2_phy_s3_q05_1}
\qn{6} A uniform disc of radius R having charge Q distributed uniformly all over its surface is placed on a smooth horizontal surface. A magnetic field B = Kxt$^2$, where K = constant, x is the distance (in metre) from the centre of the disc and t is the time (in second) is switched on perpendicular to the plane of the disc. The torque (in N-m) acting on the disc after 15 sec. (Take 2KQ = 1 S.I. unit and R = 1 metre).
\qn{7} A small ring R of mass 0.2 kg is threaded onto a fixed rod which is inclined at an angle of 30$^\circ$ to the horizontal. The ring is attached to a block B of mass 0.5 kg by means of a light string which passes over a smooth pulley P. The part RP of the string is perpendicular to the rod, and the part PB of the string is vertical (see diagram). The system is in equilibrium. Assuming the equilibrium is limiting, if the coefficient of friction ($\mu$) between R and the rod is $\mu = \dfrac{1}{\alpha - \sqrt{\beta}}$ where $\alpha$ and $\beta$ are the smallest possible integer. Fill the value of ($\alpha$ + $\beta$).\par
\qfig{p2_phy_s3_q07_1}
\qn{8} Two uniform rods of same length but different density are arranged as shown in the diagram. Centre of mass of this arrangement lies on a line ax + by = $\ell$. ($\ell$ is the length of rod) find a + b = ?\par
\qfig{p2_phy_s3_q08_1}
\partsec{2}{CHEMISTRY}
\sect{SEC-1}
\qn{1} For 0.1 M aq. solution of Na$_2$CO$_3$ at 25$^\circ$C, select the correct options (Use : H$_2$CO$_3$ : K$_{a_1}$ = 10$^{-5}$, K$_{a_2}$ = 10$^{-9}$ ) and K$_w$ (H$_2$O) = 10$^{-14}$
\begin{qoptions}
\opt{A}{pH of the solution is 11.}
\opt{B}{Equilibrium concentration of H$_2$CO$_3$ is 10$^{-9}$ M.}
\opt{C}{Equilibrium concentration of HCO$_3^-$ is 10$^{-3}$ M}
\opt{D}{Equilibrium concentration of CO$_3^{2-}$ is 0.01 M}
\end{qoptions}
\qn{2} Select the incorrect options :
\begin{qoptions}
\opt{A}{Gold sol is a positive charged sol.}
\opt{B}{Lyophobic sols can act as protective colloid for Lyophilic sols.}
\opt{C}{When FeCl$_3$ is hydrolysed with hot water, then Ferric hydroxide sol (positively charged) will be obtained.}
\opt{D}{Blood is negatively charged sol.}
\end{qoptions}
\qfullbegin
\qn{3} \qfig{p2_chem_s1_q03_1}
\qfullsplit
Select the correct statements(s) for above reaction sequence ?
\begin{qoptions}
\opt{A}{F and G are identical.}
\opt{B}{E and H are optically inactive due to external compensation.}
\opt{C}{Compound (A) and compound (B) can be differentiated by Tollen's reagent}
\opt{D}{F and G are enantiomers}
\end{qoptions}
\qfullend
\qn{4} Among the following pair of compounds, identify the pair(s) which can be differentiated by Fehling as well as Tollen's reagent ?
\begin{qoptions}
\opt{A}{\qfigi{p2_chem_s1_q04_A}}
\opt{B}{\qfigi{p2_chem_s1_q04_B}}
\opt{C}{\qfigi{p2_chem_s1_q04_C}}
\opt{D}{\qfigi{p2_chem_s1_q04_D}}
\end{qoptions}
\qfullbegin
\qn{5} Compound (A) has molecular formula C$_6$H$_7$N, undergoes following reaction sequence.\par
\qfig{p2_chem_s1_q05_1}
\qfullsplit
all are major organic product in above reaction sequence. Select the correct statement(s) for the reaction sequence.
\begin{qoptions}
\opt{A}{Order of basic strength is A $>$ D $>$ E .}
\opt{B}{Compound (A) on reaction with CHCl$_3$, Alc. KOH results formation of \qfigi{p2_chem_s1_q05_2}}
\opt{C}{Reactivity order towards aromatic electrophilic substitution is A $>$ B.}
\opt{D}{Compound (E) on reaction with SnCl$_2$, HCl followed by H$_3$O$^+$ gives \qfigi{p2_chem_s1_q05_3} as a major product.}
\end{qoptions}
\qfullend
\qn{6} Which of the following is/are true for the complex compound with IUPAC name as : Dichloridobis(ethane-1, 2-diamine)cobalt(III) sulphate
\begin{qoptions}
\opt{A}{Its formula is [Co(NH$_2$CH$_2$CH$_2$NH$_2$)$_2$Cl$_2$]SO$_4$}
\opt{B}{The co-ordination number of central metal ion is 6}
\opt{C}{The effective atomic number of central metal is 32}
\opt{D}{The compound can be categorised as a heteroleptic complex}
\end{qoptions}
\qn{7} The compound(s) that give more than one gaseous product on strong heating with Bunsen burner is/are
\begin{qoptions}
\opt{A}{Baking soda}
\opt{B}{Washing soda}
\opt{C}{Chile salt petre}
\opt{D}{Calcium oxalate}
\end{qoptions}
\qn{8} The reaction Fe + dilH$_2$SO$_4$ $\to$ ; gives the product(s)
\begin{qoptions}
\opt{A}{FeSO$_4$}
\opt{B}{Fe$_2$(SO$_4$)$_3$}
\opt{C}{H$_2$}
\opt{D}{SO$_2$}
\end{qoptions}
\sect{SEC-2}
\paraq{1}{2}For reversible Galvanic cell, thermodynamic parameters are given as :\par
$(\Delta G)_{T,P} = \Delta H - T\Delta S$\par
$(\Delta G)_{T,P} = -\,nFE_{cell}$\par
$\Delta S = nF\left[\dfrac{dE_{cell}}{dT}\right]_P$\par
$W_{usefull} = -(\Delta G)_{T,P}$\par
Answer the following questions :\par
\qn{1} For given cell Pt $|$ H$_2$ $|$ HCl $|$ AgCl $|$ Ag, variation of e.m.f. of the cell with temperature is given as E$_{cell}$ (in volt) = 0.4 $-$ 2 $\times$ 10$^{-6}$ $\left(\dfrac{T}{K}\right)^2$ (T = temperature in Kelvin) then $\Delta$G for given cell reaction at 27$^\circ$C will be $\dfrac{1}{2}$H$_2$ + AgCl $\to$ Ag + H$^+$ + Cl$^-$
\begin{qoptions}
\opt{A}{$-$ 0.44 F}
\opt{B}{$-$ 0.4 F}
\opt{C}{$-$ 0.22 F}
\opt{D}{$-$ 0.33 F}
\end{qoptions}
\qn{2} For above cell reaction, entropy change per mole of Ag at 27$^\circ$C will be (Use : 1F = 96500 C)
\begin{qoptions}
\opt{A}{$-$ 115.8 J/K}
\opt{B}{$-$ 231.6 J/K}
\opt{C}{$-$ 57.6 J/K}
\opt{D}{$-$ 463.2 J/K}
\end{qoptions}
\paraq{3}{4}The hydrolysis reactions of halides can take place at room temperature or slightly higher temperature when the vacant orbitals is/are available on the central atom and some times on the bonded atoms.\par
\qn{3} Which compounds undergo partial hydrolysis :
\begin{qoptions}
\opt{A}{BF$_3$}
\opt{B}{SiCl$_4$}
\opt{C}{SiF$_4$}
\opt{D}{SOCl$_2$}
\end{qoptions}
\qn{4} Which of the following can be hydrolysed at room temperature :
\begin{qoptions}
\opt{A}{PCl$_3$}
\opt{B}{PCl$_5$}
\opt{C}{NF$_3$}
\opt{D}{NCl$_3$}
\end{qoptions}
\sect{SEC-3}
\qn{1} A gaseous mixture of H$_2$ and He in the mass ratio 1 : 4 is successively effused through a series of containers connected through narrow tubes of negligible volume and fitted with valves. The valves are successively opened for a very small time interval, dt. Find the minimum number of total containers required to obtain a molar ratio 8 : 1 for H$_2$ and He after the effusion process.
\qn{2} One mole of an diatomic ideal gas in an insulated cylinder fitted with a massless, frictionless movable piston is expanded gradually from 1 L to 32 L. If the initial temperature of gas is 1000 K then find magnitude of change in Gibb's Free energy in kcal.\par
(Molar entropy of gas = 3 cal mol$^{-1}$/k$^{-1}$)\par
(Use, R = 2 cal mol$^{-1}$ k$^{-1}$)\par
(Assume molar entropy of a gas is constant)
\qn{3} The activation energy of a reaction is given as E$_a$ (cal mol$^{-1}$) = 0.01T$^2$ + $\dfrac{20}{7}$T where T is temperature in kelvin. Find value of $\ell$n$\dfrac{k_2}{k_1}$ where k$_1$ and k$_2$ are rate constants at temperature 200 K and 400 K respectively. (Use 2 = 0.7 and R = 2 cal mol$^{-1}$ k$^{-1}$)
\qn{4} A polyhydroxy compound having molecular mass 150, on reacting the compound with CH$_3$COCl, the precipitate obtained has molecular mass of 318. Calculate number of OH group present in compound :
\qn{5} Examine the given structure and select those that satisfy given conditions\par
\qfig{p2_chem_s3_q05_1}
Find out no. of compounds which give SN$^2$ reaction with in CH$_3\overset{\ominus}{\mathrm{S}}\,\overset{\oplus}{\mathrm{Na}}$ in DMF faster with respect to CH$_3$CH$_2$Br
\qn{6} (a) Total number of stereoisomers of given compound are x\par
\qfig{p2_chem_s3_q06_1}
(b) Number of optically active stereoisomers of 1,2,3,4,5,6-hexachlorocyclohexane is y.\par
Then x $-$ y is
\qn{7} In the list of compounds given below, identify the number of compounds that have d$\pi$ $-$ p$\pi$ bonding in the molecular structure with sp$^2$ hybridization of the underlined atom.\par
$\underline{\mathrm{C}}$Cl$_4$, $\underline{\mathrm{P}}$OCl$_3$, $\underline{\mathrm{S}}$O$_2$, $\underline{\mathrm{S}}$O$_3$, H$_2\underline{\mathrm{S}}$O$_4$, $\underline{\mathrm{S}}$OCl$_2$, $\underline{\mathrm{S}}$O$_2$Cl$_2$, $\underline{\mathrm{B}}$F$_3$, $\underline{\mathrm{B}}$Cl$_3$, $\underline{\mathrm{N}}$(SiH$_3$)$_3$
\qn{8} How many of the following is/are yellow and water insoluble substances ?\par
Ag$_2$CrO$_4$, PbCrO$_4$, Na$_2$CrO$_4$, BaCrO$_4$, CrO$_2$Cl$_2$, (NH$_4$)$_2$Cr$_2$O$_7$, (NH$_4$)$_3$PO$_4$.12MoO$_3$, K$_3$[Co(NO$_2$)$_6$], AgI, PbI$_2$
\partsec{3}{MATHEMATICS}
\sect{SEC-1}
\qn{1} If $f(x) = 1 - (\sin 2x)$\par{\footnotesize$\sqrt[4]{2\sqrt{\tan^2 x((\tan^2 x) + 6) + 9} + \tan^4 x - 5}$}\par where x $\in \left(\dfrac{3\pi}{2}, 2\pi\right]$
\begin{qoptions}
\opt{A}{$f(x) = 1 - 2\sin x$}
\opt{B}{$f(x) = 1 + 2\cos x$}
\opt{C}{$f(x)$ can take only 2 integral value}
\opt{D}{$f(x)$ can take only 3 integral value}
\end{qoptions}
\qn{2} x + y = 2 and x $-$ y = 2 are tangents on a parabola at (1, 1) and (4, 2) respectively. Which of the following is/are correct
\begin{qoptions}
\opt{A}{Equation of directrix is x + 3y = 2}
\opt{B}{Equation of axis is 3x $-$ y = 5}
\opt{C}{Focus of the parabola is at $\left(\dfrac{8}{5}, \dfrac{6}{5}\right)$}
\opt{D}{Vertex of the parabola is at $\left(\dfrac{33}{20}, \dfrac{13}{20}\right)$}
\end{qoptions}
\qn{3} Probability of n heads in 2n tosses of a fair coin can be given by
\begin{qoptions}
\opt{A}{$\displaystyle\prod_{r=1}^{n}\left(\dfrac{2r - 1}{2r}\right)$}
\opt{B}{$\displaystyle\prod_{r=1}^{n}\left(\dfrac{n + r}{2r}\right)$}
\opt{C}{$\displaystyle\sum_{r=0}^{n}\left(\dfrac{{}^nC_r}{2^n}\right)^2$}
\opt{D}{$\dfrac{\displaystyle\sum_{r=0}^{n}\left({}^nC_r\right)^2}{\left(\displaystyle\sum_{r=0}^{2n}{}^{2n}C_r\right)}$}
\end{qoptions}
\qn{4} P(n) denotes the number of natural number less than n, which are relatively prime with n, then
\begin{qoptions}
\opt{A}{P $\left(2^2.5\right) = 2^3$}
\opt{B}{P $\left(2^3\right) = 2^2$}
\opt{C}{P $\left(2^2.3\right) = 2^2$}
\opt{D}{P $\left(2.3^2\right) = 2^2.3$}
\end{qoptions}
\qn{5} If the line y = 2 $-$ x is tangent to the circle `s' at a point P(1, 1) and circle `s' is orthogonal to the circle x$^2$ + y$^2$ + 2x + 2y $-$ 2 = 0 then which of the following statement is true about circle `s'.
\begin{qoptions}
\opt{A}{The equation of the circle `s' is x$^2$ + y$^2$ $-$ x $-$ y = 0}
\opt{B}{The length of the tangent from (2, 2) to the circle `s' is `2'}
\opt{C}{The equation of the director circle of circle `s' is x$^2$ + y$^2$ $-$ 2x $-$ 2y + 2 = 0}
\opt{D}{The length of the common chord of given circle and circle `s' will be $\dfrac{\sqrt{10}}{3}$}
\end{qoptions}
\qn{6} If I$_m$ = $\displaystyle\int_0^{m\pi} x\,|\sin x|\,e^{\cos^4 x}\,dx$ and J$_n$ = $\displaystyle\int_0^{n\pi} |\cos x|\,e^{\cos^4 x}\,dx$, where m, n $\in$ I then\par
(Note : [k] denote GIF)
\begin{qoptions}
\opt{A}{$\left[\dfrac{I_1}{J_1}\right] = 0$}
\opt{B}{$\left[\dfrac{I_2}{J_2}\right] = 3$}
\opt{C}{$\left[\dfrac{I_2}{J_1}\right] = 0$}
\opt{D}{$\left[\dfrac{I_3}{J_3}\right] = 4$}
\end{qoptions}
\qn{7} Let f (x) = cos x sin 2x, then
\begin{qoptions}
\opt{A}{min f (x) $> -\dfrac{7}{9}\,\forall x \in (-\pi, \pi)$}
\opt{B}{min f (x) $> -\dfrac{9}{7}\,\forall x \in (-\pi, \pi)$}
\opt{C}{min f (x) $> -\dfrac{1}{9}\,\forall x \in [-\pi, \pi]$}
\opt{D}{min f (x) $> -\dfrac{2}{9}\,\forall x \in [-\pi, \pi]$}
\end{qoptions}
\qn{8} Area bounded by $\sqrt{x + y + 1} < \sqrt{y - x + 2}$ \& y $<$ 0 is less than
\begin{qoptions}
\opt{A}{$\dfrac{9}{16}$}
\opt{B}{$\dfrac{9}{5}$}
\opt{C}{$\dfrac{5}{4}$}
\opt{D}{$\dfrac{16}{5}$}
\end{qoptions}
\sect{SEC-2}
\paraq{1}{2}A square matrix A = [a$_{ij}$] is called symmetric matrix if a$_{ij}$ = a$_{ji}$ $\forall$ i, j and a square matrix A = [a$_{ij}$] is called skew -symmetric matrix if a$_{ij}$ = $-$a$_{ji}$ $\forall$ i, j.\par
\qn{1} If A and B are symmetric matrices of order 3 then, Which of the following statements are ALWAYS TRUE.\par
I) A$^{2n+1}$ is a skew-symmetric matrix for n $\in$ N\par
II) AB + BA is a symmetric matrix,\par
III) AB $-$ BA is a skew-symmetric matrix\par
IV) Det(B) is non zero.
\begin{qoptions}
\opt{A}{I only}
\opt{B}{II \& III}
\opt{C}{I, II \& III}
\opt{D}{I, II , III \& IV}
\end{qoptions}
\qn{2} A = $\begin{bmatrix} 0 & 2 & -3\\ -2 & 0 & 5\\ 3 & -5 & 0 \end{bmatrix}$, X = $\begin{bmatrix} x\\ y\\ z \end{bmatrix}$, B = $\begin{bmatrix} 0\\ 0\\ 0 \end{bmatrix}$ Then AX = B, will have
\begin{qoptions}
\opt{A}{Exactly one solution for X}
\opt{B}{No solution for X}
\opt{C}{Multiple solutions for X}
\opt{D}{None of these}
\end{qoptions}
\paraq{3}{4}Linked comprehensive type passage:\par
Equation of line L = 0 is given by 2x $-$ y + z $-$ 2 = 0 = x + 2y $-$ z $-$ 3 , then\par
\qn{3} The equation of the plane containing the line L = 0 and passing through the point (3, 2, 1) is ax + by + cz + 1 = 0, then a $-$ b + c is divisible by
\begin{qoptions}
\opt{A}{3}
\opt{B}{4}
\opt{C}{2}
\opt{D}{6}
\end{qoptions}
\qn{4} Vector equation of the plane which passes through the point ($-$1, 3, 2) and is perpendicular to line L = 0 is
\begin{qoptions}
\opt{A}{$\vec{r}\cdot(\hat{i} - 3\hat{j} - 5\hat{k}) + 20 = 0$}
\opt{B}{$\vec{r}\cdot(\hat{i} - 3\hat{j} + 5\hat{k}) = 0$}
\opt{C}{$\vec{r}\cdot(\hat{i} + 3\hat{j} + 5\hat{k}) - 18 = 0$}
\opt{D}{$\vec{r}\cdot(-\hat{i} + 3\hat{j} - 5\hat{k}) + 3 = 0$}
\end{qoptions}
\sect{SEC-3}
\qn{1} A tangent is drawn to the curve, $\dfrac{x^2}{16} + \dfrac{y^2}{9} = 1$ at the point P meeting the co-ordinate axis in T and t. If perpendicular from origin O to the about tangent meet at Y, then value of the product (Tt)(PY) is equal to
\qn{2} A continuous function f : $\left[0, \dfrac{3}{4}\right] \to$ R satisfies the equation $f(x) = (1 + x^2)\left(1 + \displaystyle\int_0^x \dfrac{f^2(t)}{1 + t^2}\,dt\right)$, then $\left[f\left(\dfrac{1}{2}\right)\right]$ is\par
[Note : [K] denotes greatest integer less than or equal to K.]
\qn{3} The quadratic polynomial P(x) satisfy the equation P(x) $-$ P(x + 1) = 2x $-$ 1 $\forall$ x $\in$ R. If the greatest value of P(x) in ($-\infty$, 0] is 2, then value of P($-$1) + P(1) is equal to
\qn{4} S$_n$ = $\displaystyle\lim_{x \to \infty}\bigl(\bigl($ $\left(x + \dfrac{1}{\sqrt{4n^2}}\right)$ $\left(x + \dfrac{1}{\sqrt{4n^2 - 1^\alpha}}\right)$ $\left(x + \dfrac{1}{\sqrt{4n^2 - 2^\alpha}}\right)$ $\cdots$ $\left(x + \dfrac{1}{\sqrt{4n^2 - (n-1)^\alpha}}\right)$ $\bigr)^{1/n} - x\bigr)$, n$\in$N\par
and for $\alpha$ = 1, L = $\displaystyle\lim_{n \to \infty}$ (n. S$_n$) and\par
for $\alpha$ = 2, M = $\displaystyle\lim_{n \to \infty}$ (n. S$_n$)\par
then value of [L + M] is equal to\par
[Note : [y] denotes the greatest integer function of y.]
\qn{5} Let a real valued function $f(x) = x^4 + 4x^3 - 8x^2 + P$, P $\in$ R . If $\lambda$ is the least value of P for which $g(x) = |f(x)|$ is differentiable $\forall$ x $\in$ R , then value of $\dfrac{\lambda}{16}$ is equal to
\qn{6} Let A and B be two sets such that n(A) = 4 and n(B) = 3. If (x, 1), (y, 2), (z, 1) and (2, 3) are in A $\times$ B, where x, y, z are distinct elements. If a function $f$ : A $\to$ B defined \& number of onto function from set A to set B is $\lambda$, then value of $\dfrac{\lambda}{6}$ is equal to
\qn{7} Let 3 vector $\vec{\ell}$, $\vec{m}$ \& $\vec{n}$ are such that $\left|\vec{\ell}\right| = 3$, $\left|\vec{m}\right| = \dfrac{1}{2}$ \& $\vec{\ell}\cdot\vec{m} = \vec{m}\cdot\vec{n} = \vec{n}\cdot\vec{\ell} = 0$. If $\left[\vec{\ell}\times\left(\vec{m}\times\vec{\ell}\right)\ \ \ \vec{n}\times\vec{m}\ \ \ \vec{n}\right] = 36$, then value of $\left|\vec{n}\right|$ is equal to
\qn{8} Let a circle with centre O and radius k. 3 points P, Q, R are chosen on a circle such that if length of chord PQ, QR, RP are x, y \& z respectively, x$^2$ + y$^2$ + z$^2$ = 5, $\angle$QOR = 60$^\circ$ and Area of $\Delta$PQR is $\dfrac{\sqrt{3}}{4}$ unit, then the value of radius of circle k will be equal to
\end{multicols}
\end{document}
